
\AddSymbolsD
\pagecolor{Goldenrod!20!Gray}
\color{white}
\quad\quad \lettrine[,lraise=0.2,lhang=0.5]{\textbf{教}}{授}和铃木和子走在晴朗的夜空之下。繁星点点，勾勒出夜空中少有的宁静。

铃木和子抬头望着星空，任由微弱的星光抵达自己的瞳孔；她的手里则仿佛不经意间多拿了一张伊藤佐一留下的草稿纸。

教授则不顾天空之屏，低头喃喃自语着。

\twkkai{第163颗星……上面会有什么呢？}

铃木和子听不懂教授的喃喃自语，忐忑地向前走着。

\twkkai{如果伊藤佐一在身边，该多好啊……}

这次，铃木和子没有留意到自己脑海中多余的情愫，任由其扩散弥漫着。

\twkkai{我该问教授什么问题好呢？}她沉思着。

转眼间，他们走到了一间茶舍。教授率先发话。

\twkkai{我们坐下品一下茶吧？最近我的思绪有点乱……品完茶我便要先回去了——你有什么问题可以先问我；}教授平静地看着铃木和子。

铃木和子拿出了伊藤佐一画的那张图。教授看着字迹，推了推眼镜，嘴里继续呢喃着什么。

\twkkai{这小子……没想到啊；}教授的嘴角弯了起来，又很快地被他压了下去。

教授很快便和铃木和子讲解起来。讲解完毕后，教授突然问了一句。

\twkkai{你是想我跟你倾诉好，还是不倾诉好呢？}他看着铃木和子澄澈的双眼，说着。

铃木和子不解地看着教授；教授没有让空气安静太久。

\twkkai{没关系；如果你能给我一点线索，那最好——这一切都和一个数字有关；}教授在茶桌上比划着。

铃木和子突然有不好的预感。她早已知悉自己和奖章的不可明说的关系；她尚且不知道伊藤佐一是否理解自己的处境——她渴望他理解，但不敢强求他理解；此刻，听到教授再次讲一些奇形怪状的数学，她感到的不是新奇，而是焦虑和恐慌——或许还有一点虚幻的疼痛。

\twkkai{或许就是……真理的代价？}

她想着，然后连忙摇了摇头。

\twkkai{教授，我不太懂……}

\twkkai{没事，没事——不用放在心上；它本应该困扰的是我，而不是你，铃木同学；}教授低着头，摆摆手。

铃木和子愣了一下。那一刻，她突然间感到一阵说不清的心疼和愧疚——教授虽然在求救，但同时也似乎在奋力阻止着自己把她卷入什么。

她也并没有想明白一个问题：一个充满幽默感的人物，怎么会为一个小数字困扰呢？

而教授终于又发话了。

\twkkai{这个数字一出现在我的脑海，我就知道有什么不妥的事要发生了；}他捂着脸说。

铃木和子无奈又无助地左右张望着，突然从茶舍的窗外瞥见了什么。她拍了拍教授。

\twkkai{教授，看……}

教授早已被自己的思绪和话语弄得疲惫，只能无力地看向窗外；结果一扭头，他立即颤栗起来：

\begin{figure}[htbp]
  \centering
  \scalebox{1}{ 
  \begin{tabular}{ccc}
    \includegraphics[width=0.3\textwidth]{y_2_x_3_1 (3).pdf} &
    \includegraphics[width=0.3\textwidth]{y_2_x_3_x_1 (2).pdf} &
    \includegraphics[width=0.3\textwidth]{y_2_x_3_2x_1 (1).pdf} \\
    \includegraphics[width=0.3\textwidth]{y_2_x_3_3x_2 (1).pdf} &
    \includegraphics[width=0.3\textwidth]{y_2_x_3_12x_16 (1).pdf} &
    \includegraphics[width=0.3\textwidth]{y_2_x_3_27x_54 (1).pdf} \\
    \includegraphics[width=0.3\textwidth]{y_2_x_3_8 (1).pdf} &
    \includegraphics[width=0.3\textwidth]{压轴.pdf} &
    \includegraphics[width=0.3\textwidth]{Yahoo.pdf} \\
  \end{tabular}
  }
\end{figure}
\clearpage

原来，窗外的繁星已经依次列成了几幅奇异曲线的图案；这让两人都惊异不已。

\twkkai{这就很奇怪了……铃木同学，你可能不知道——这些星星，就是和163关联很大的……；}

教授仿佛在给铃木和子耐心地讲授，实则声音中透露出一阵让铃木和子也逐渐感到不安的颤抖。

\twkkai{其实，这个数字为什么会困扰你呢？}铃木和子略显紧张地说着，刘海因为汗水而开始紧贴着她的额头，弄得她极为不舒服。

\twkkai{铃木同学，我该怎么和你说呢？你手上的纸，也和我的命运息息相关——在此，我不卖关子了，和你讲清楚吧……}

他的手颤抖地写着，可字迹依旧是那么平稳。
\begin{equation*}
    \frac{1}{\pi}
    = \frac{12}{\sqrt{640320^{3}}}
      \sum_{n=0}^{\infty}
      (-1)^{n}
      \frac{(6n)!}{(3n)!\,(n!)^{3}}
      \cdot
      \frac{13591409 + 545140134\,n}{640320^{3n}}
\end{equation*}

\twkkai{铃木同学，这个级数，由我告诉你吧——就是用来求圆周率的；却让我也得到了无穷的苦役……}

\twkkai{为什么呢？不就是……？}铃木和子问道。

\twkkai{可不是嘛——就是一个圆周率；}教授自嘲着，又倒吸了一口冷气。

\twkkai{但它已经让我做了一个月噩梦了。每晚，总会梦到一位老者，厉声痛斥我，让我在月球上漫步，以精确测量月球的半径——他也谓之“守门”。有时走到太阳晒到的岩石表面，脚底烫得发麻，却发现走快走慢也无济于事；}教授仿佛有点呼吸困难，瞥着窗外的那些繁星，始终无法平复。

铃木和子替教授担心着；可是她总觉得自己是“五十步忧百步”。她也有属于自己的噩梦；她不知道她心上的那个人有没有。她多么希望噩梦自己也有噩梦。

教授又开口了。

\twkkai{铃木同学，我想说——真正蹊跷的是，这个级数不是我乱编的；它是真实的、正确的，而且和163的渊源可深了。可以说，老者给我下的魔咒——就是163。他不打算让人听懂；}教授的声音越来越大了。

\twkkai{说什么胡话呢……；}茶舍里的人们纷纷暗想着。只有铃木和子不这么想——她知道。

繁星不打算消散。教授却也不打算退缩。他用大拇指悄悄地指了指群星，示意着铃木和子。

\twkkai{至于这些——我是不打算理会的；但有必要和你说清楚。}

一阵眩晕感向铃木和子袭来。她虽然试图理解教授；但教授的话似乎太重——重到她以为是道德绑架。

\twkkai{不要看这些椭圆曲线长得花枝招展——其实它们内部早便有了自己作为天庭发号施令者的身份证……}

在人们眼中，教授像喝醉了茶的怪人；几位临近的客人连忙起身找老板结了账。

\twkkai{它们的身份证便是\(j(\tau)\)。我和你说啊，铃木同学——找到它的值，你就抓住了这些曲线的把柄了！}

教授咳嗽了几声。铃木和子听到他的咳嗽，反而终于安下了心——她知道。

什么将要回归正常了。

\twkkai{
实际上啊——$j(\tau)$（也就是$j$-不变量）是模函数中最重要的一个，长这样：
\[
j(\tau) = \frac{1}{q} + 744 + 196884q + 21493760q^2 + \cdots
\]
你可以先不用管那些系数是怎么来的——记住它长这样就可以了。
}\index{jbubianliang@$j$-不变量 $j(\tau)$}
\index{jbubianliang@$j$-不变量 $j(\tau)$!傅里叶展开 $q^{-1}+744+196884q+\cdots$}


铃木和子没有吭声。她看着那些系数——196884,21493760……满脸惊奇——尽管她表现得很平静。她听着教授说“记住”，可是怎么记得住呢？

\twkkai{简直比学生证号码还难记呢；}她默想着。

人们的笑声稀疏了——或许是在什么中走散了吧。而教授的手依旧颤巍巍地写着。

\twkkai{
这时我们要引入一个级数——它叫艾森斯坦级数：
\[
G_k(\tau) = \sum_{(m,n) \neq (0,0)} \frac{1}{(m + n\tau)^k}
\]

取 $k=4$ 和 $k=6$——然后，归一化便得到
\[
E_4(\tau) = 1 + 240\sum_{n=1}^\infty \sigma_3(n)q^n,
\quad
E_6(\tau) = 1 - 504\sum_{n=1}^\infty \sigma_5(n)q^n
\]
}\index{moxing@模形式!艾森斯坦级数 $G_k(\tau)$、$E_4(\tau)$、$E_6(\tau)$}\index{moxing@模形式!归一化艾森斯坦系数（240、504 与伯努利数）}


教授在其中一个可怜的符号上圈圈点点着，不给铃木和子听到他详细解释的机会。可是她多想帮教授、也帮自己弄清楚——那蹊跷的240和504究竟是怎么来的啊。

月亮在窗的边缘躲着，赶走了星星。那些天空中的曲线消散了。铃木和子松了一口气，却多了股莫名的哀伤——仿佛她要为那星星哀悼。

\twkkai{
其中：
\[
\sigma_k(n) = \sum_{d|n} d^k.
\]

铃木同学，这个$\sigma_k(n)$你应该记得吧？}\index{chushuhe@除数和函数$\sigma_k(n)=\sum_{d\mid n}d^k$}

铃木和子没有反应过来；她的眼神迷离。教授没有抬头也感受到了——从那比平常漫长的停顿中。

\twkkai{
不用紧张——不是来考你的，只是确认一下罢了；}教授推了推眼镜，看着铃木和子。

此时，两个人都很紧张——但紧张的对象不尽相同。

\twkkai{好了，好了——接下来是压轴戏！}教授没有理会自己，也没有理会别人——没有什么好理会的。值得理会的只有他脑海中的苦役。

\twkkai{

再定义判别式函数（拉马努金$\Delta$函数）：
\[
\Delta(\tau)
=
\frac{E_4(\tau)^3 - E_6(\tau)^2}{1728}
=
q\prod_{n=1}^\infty (1-q^n)^{24}
\]

它是权为 $12$ 的尖点形式，在 $\mathbb{H}$ 上处处不为零。}\index{moxing@模形式!$\Delta(\tau)=(E_4^3-E_6^2)/1728$（权$12$尖点形式）}\index{moxing@模形式!$\Delta(\tau)$ 在 $\mathbb{H}$ 上处处不为零}\index{lamanujin@拉马努金（S.\,A.\ Ramanujan）!拉马努金 $\Delta$ 函数 $\Delta(\tau)$}

\twkkai{处处不为零——是什么意思呢？还有，刚才的240和504，是怎么……？}铃木和子把椅子往前移了移。

\twkkai{啊……240和504——又是两个缠人的小妖精罢了。你想知道，我就大概和你点一点吧。它们都是这个形式的数：\(-\frac{2k}{B_{2k}}\)——其中\(B_{2k}\)是伯努利数，又和黎曼\(\zeta\)函数有关。只不过那黎曼，简直比狄利克雷还要难缠。要解释的话，还是改天吧。}\index{bonulishu@伯努利数$B_{2k}$}

教授无奈地摇摇头，用手扶着脑袋。铃木和子能感受到他心中的苦根有多深。\twkkai{我可以帮你验证——它们实际上是黎曼\(\zeta\)函数在偶数点的值\index{li@黎曼（B. Riemann）!$\zeta(s)$ 在偶数点的值}。代入：$k=4$：$B_4 = -\dfrac{1}{30}$，那么 $-\dfrac{2\cdot4}{B_4} = -\dfrac{8}{-1/30} = 240$；$k=6$：$B_6 = \dfrac{1}{42}$，于是 $-\dfrac{2\cdot6}{B_6} = -\dfrac{12}{1/42} = -504$……}

铃木和子看了看教授布满皱纹的手，不敢说什么。

\twkkai{
而伯努利数和\(\zeta\)函数的关系是：
$$\zeta(2k) = \frac{(-1)^{k+1}(2\pi)^{2k}}{2\cdot(2k)!}B_{2k}.$$所以240和504，本质上是$\zeta(4)=\pi^4/90$ 和 $\zeta(6)=\pi^6/945$这两个家伙派来对付我的。}\index{bonulishu@伯努利数$B_{2k}$!与 $\zeta(2k)$ 的关系}

铃木和子看不懂——但她也舍不得将视线离开那些神谕般的符号。\twkkai{原来，伊藤同学每天面对的就是这些吗……；}她不敢相信这一切，发生在自己身上的同时，也发生在了他身上。

\twkkai{至于你说的处处不为零——这对于那些星星组成的队列是必须的；不然它们会塌……；}教授朝外面天空中不复存在的星星队列努了努嘴。

\twkkai{会塌？}

\twkkai{是的；像人一样，需要有价值观支撑他活着。}

教授终于抬起了头，郑重地看着铃木和子。她再次感受到一种重量——不过这次不一样。\twkkai{或许是——积极的？}她对自己所思所想感到莫名其妙。

\twkkai{
好了，那么——$j(\tau)$——那些星星的把柄，就被定义为：
\[
\boxed{
j(\tau)
=
\frac{E_4(\tau)^3}{\Delta(\tau)}
}
\]}\index{jbubianliang@$j$-不变量 $j(\tau)$!定义 $j(\tau)=E_4(\tau)^3/\Delta(\tau)$}

教授默默地把那条定义式框了起来——也仿佛把自己囚禁了进去。他继续默默地写着。

\twkkai{它不会随着模群的变换而改变——}他接着说，拿起笔写道：
\[
\tau \mapsto \frac{a\tau+b}{c\tau+d}
\]

\twkkai{无论怎么变换——}
\[
j\!\left(\frac{a\tau+b}{c\tau+d}\right) = j(\tau)
\]

\twkkai{它都如如不动——一个权为0的模函数，完全不会变。}

铃木和子仿佛听出了教授言语中的沮丧；她便又想起了刚才的 $\Delta$——处处不为零；而现在 $j(\tau)$ 完全不变。\twkkai{两者有何联系呢？}她默默地想，但不敢说出来。

而教授早已用笔描摹出一份图表了：

\begin{table}[h]
\centering
\renewcommand{\arraystretch}{1.2}
\begin{tabular}{cc}
\toprule
$\tau$ & $j(\tau)$ \\
\midrule
$i$ & $1728$ \\
$e^{2\pi i/3}$ & $0$ \\
$e^{2\pi i/4}$ & $1728$ \\
$\dfrac{1+\sqrt{-163}}{2}$ & $-640320^3$ \\
\bottomrule
\end{tabular}
\end{table}
\index{jbubianliang@$j$-不变量 $j(\tau)$!$j\left(\frac{1+\sqrt{-163}}{2}\right)=-640320^3$}

写毕，教授突然侧着俯身拿起一壶水，往茶杯里倒。\twkkai{先不要理这份表；}他平静地说。

\twkkai{数数茶吧。}

茶水汩汩地响着。茶杯口盖着一只不锈钢漏斗；漏斗底部正滴着茶。教授目不转睛地盯着漏斗，看着一滴滴茶在茶杯肚子里激起浪花、不断翻滚。

教授的心情也随之翻滚着。

1，2，3，……

7，11，19，……

43，67，……

163……

茶滴停下来了。而他又想起自己还未完成的苦役。

\twkkai{刚才我数的，就是Heegner数——你有听到吗？}\index{heegner@Heegner 数}教授轻声地说——是那么轻，以至于铃木和子是真的差点听不到了。她也轻轻点了点头。

\twkkai{你看啊，这份表里有两个数比较蹊跷——一个是1728，即\(12^3\)，它是\(\Delta\)的分母；另一个便是……}

教授停顿了一下，又突然意识到自己是在铃木和子面前，便又被迫继续说着。

\twkkai{……另一个便是\(-640320^3\)这个数了。它怎么来的呢——又和最大的Heegner数163有关了。至此，你应该能感受到我一半的苦恼了吧；}\index{heegner@Heegner 数!最大 Heegner 数 $163$}教授摇了摇头。

此时，茶舍已经只剩寥寥几人了。窗外的风忽然闯了进来，撩动着教授的头发和铃木和子的刘海；铃木和子下意识地挡了挡。教授并无心留意这一切。

\twkkai{我写下来吧：
\[
\tau_{163}
=
\frac{1+\sqrt{-163}}2.
\]
对应：
\[
j\!\left(\frac{1+\sqrt{-163}}2\right)
=
-262537412640768000
=
-640320^3.
\]}

\twkkai{
由于类数为 $1$：
\[
j(\tau)\in\mathbb Z.
\]}\index{leishu@类数（Class number）}\index{leishu@类数（Class number）!类数为 $1$ 时 $j(\tau)\in\mathbb{Z}$}

\twkkai{什么是类数？}铃木和子知道不应该在教授疲惫的时候问话，但她还是脱口而出了。

教授看着铃木和子，又看回草稿纸，意识到什么。

\twkkai{这个下次再讲吧——我们会再品一次茶的……}

铃木和子没有听懂教授的言外之意。

\twkkai{走吧，时间不早了；}教授率先站起了身。铃木和子把最后一口茶细细地品完后，也站了起来，并拍了拍裙子。

\twkkai{哦，对了，这张纸留给你吧……我有天晚上睡不着写的，我自己也不想留着——不想让它再给我困扰；}教授从衬衫口袋里拿出一张卡纸。

铃木和子看着卡纸，呼吸急促了起来——但并不是因为紧张。

\clearpage
\myblankpage
\includepdf[rotate=270, scale=0.9]{163.pdf}

天气凉了，但那张纸的余温仍存。

教授早已离开了。

而铃木和子携着教授的心意，心怀感激地离开了茶舍。







